%%%%%%%%%%%%%%%%%%%%%%%%%%%%%%%%%%%%%%%%%%%%%%%%%%%%%%%%%%%%%%%%%%%%%%
% How to use writeLaTeX: 
%
% Este arquivo foi gerado para descrever a metodologia de classificação
% de grafos com Redes Neurais sem Peso (WiSARD), conforme os scripts
% em Python fornecidos.
%
%%%%%%%%%%%%%%%%%%%%%%%%%%%%%%%%%%%%%%%%%%%%%%%%%%%%%%%%%%%%%%%%%%%%%%

\documentclass[12pt]{article}

\usepackage{sbc-template}
\usepackage{graphicx,url}
\usepackage[utf8]{inputenc}
\usepackage{amsmath}
\usepackage{amssymb}
\usepackage{booktabs}
\usepackage{hyperref}
     
\sloppy

\title{Uma Pipeline Automatizada para Representação Binária de Grafos e Classificação com WiSARD}

\author{Hélio H. Souza\inst{1}, Diego N. Brandão\inst{2}}

\address{Centro Federal de Educação Tecnológica Celso Suckow da Fonseca (CEFET/RJ)\\
  Rio de Janeiro -- RJ -- Brazil
\nextinstitute
  Programa de Pós-graduação em Ciência da Computação (PPCIC) -- CEFET/RJ\\
  Rio de Janeiro -- RJ -- Brazil
  \email{helio.souza@aluno.cefet-rj.br, diego.brandao@cefet-rj.br}
}

\begin{document} 

\maketitle

\begin{abstract}
This paper presents a complete and automated pipeline for graph classification using Weightless Neural Networks (WNNs), focusing on the WiSARD classifier. The proposed methodology, named M+Q, consists of two main stages: (M) extraction of a rich set of topological attributes based on centrality measures; and (Q) quantization of these attributes into binary vectors. The quantization stage employs the KernelCanvas++ technique for continuous attributes and a histogram-thermometer encoding for categorical attributes. A key differentiator is the automatic bit-budget allocation, which optimizes the distribution of information between the graph's structural and non-structural attributes. We evaluate the pipeline on benchmark datasets, demonstrating its effectiveness and the importance of hyperparameter optimization in generating discriminative binary representations.
\end{abstract}
     
\begin{resumo} 
Este artigo apresenta uma pipeline completa e automatizada para a classificação de grafos utilizando Redes Neurais sem Peso (WNNs), com foco no classificador WiSARD. A metodologia proposta, denominada M+Q, consiste em duas etapas principais: (M) extração de um rico conjunto de atributos topológicos baseados em medidas de centralidade; e (Q) quantização desses atributos em vetores binários. A etapa de quantização emprega a técnica KernelCanvas++ para atributos contínuos e uma codificação de histograma com termômetro para atributos categóricos. Um dos principais diferenciais é a alocação automática de orçamento de bits, que otimiza a distribuição de informações entre os atributos estruturais e não-estruturais do grafo. Avaliamos a pipeline em datasets de benchmark, demonstrando sua eficácia e a importância da otimização de hiperparâmetros na geração de representações binárias discriminativas.
\end{resumo}


\section{Introdução}

A classificação de grafos é uma tarefa fundamental em aprendizado de máquina, com aplicações em bioinformática, análise de redes sociais e visão computacional. O desafio central reside na conversão da estrutura topológica complexa e de tamanho variável de um grafo em um vetor de características de tamanho fixo, que possa ser processado por algoritmos de classificação convencionais.

Redes Neurais sem Peso (WNNs), como a WiSARD (\textit{Wilkie, Stonham, and Aleksander's Recognition Device}) \cite{aleksander1984}, oferecem uma alternativa computacionalmente eficiente às redes neurais tradicionais. Elas operam diretamente sobre dados binários, utilizando tuplas de bits de entrada como endereços para acessar memórias de acesso aleatório (RAMs) que armazenam os padrões aprendidos. Essa arquitetura as torna extremamente rápidas em treinamento e inferência, mas impõe um pré-requisito estrito: os dados de entrada devem ser binarizados.

A conversão de dados estruturados, como grafos, em representações binárias eficazes não é trivial. Uma binarização inadequada pode levar à perda de informações discriminativas, degradando o desempenho do classificador. Diante disso, este trabalho apresenta uma pipeline completa e flexível, denominada M+Q, projetada para transformar grafos em vetores binários otimizados para WNNs. A pipeline consiste em duas etapas: (M) extração de métricas e (Q) quantização.

A principal contribuição deste artigo é uma metodologia de ponta a ponta que não apenas extrai e binariza atributos de grafos, mas também incorpora um mecanismo de alocação automática de orçamento de bits. Esse mecanismo otimiza a distribuição dos recursos de representação (bits) entre diferentes tipos de atributos, automatizando uma etapa crítica de ajuste de hiperparâmetros e melhorando o poder discriminativo dos vetores binários gerados.

\section{Metodologia Proposta: A Pipeline M+Q}

A pipeline M+Q é um processo de duas fases que transforma uma coleção de grafos em um conjunto de vetores binários e seus respectivos rótulos, pronto para ser usado no treinamento e teste de um classificador WiSARD. A Figura~\ref{fig:pipeline} ilustra o fluxo de trabalho.

\begin{figure}[ht]
\centering
% \includegraphics[width=0.9\textwidth]{pipeline_diagram.png} % Placeholder for a diagram
\caption{Visão geral da pipeline M+Q para classificação de grafos. (Diagrama a ser inserido)}
\label{fig:pipeline}
\end{figure}

\subsection{Etapa M: Extração de Atributos (Discretização Topológica)}

A primeira etapa da pipeline, denominada M (Métricas), é responsável por traduzir a estrutura topológica de cada grafo em um vetor de características numéricas e categóricas. Este processo, implementado no módulo \texttt{reading\_graphs.py}, pode ser entendido como uma forma de \textit{discretização topológica}, onde a complexa e contínua natureza das conexões de um grafo é mapeada para um conjunto finito e tabular de atributos.

O processo é executado para cada grafo do dataset (e.g., no formato TUDataset) e consiste nos seguintes passos:
\begin{enumerate}
    \item \textbf{Cálculo de Atributos Estruturais:} Para cada nó do grafo, um conjunto rico de métricas de centralidade é calculado. Estas métricas, implementadas de forma robusta em \texttt{centralities.py} para lidar com problemas de convergência ou grafos desconexos, quantificam a importância estrutural de cada nó sob diferentes perspectivas. As métricas incluem:
    \begin{itemize}
        \item \textbf{Locais:} Centralidade de Grau (\textit{Degree}).
        \item \textbf{Globais:} Centralidade de Proximidade (\textit{Closeness}) e Intermediação (\textit{Betweenness}).
        \item \textbf{Espectrais:} PageRank e Centralidade de Autovetor (\textit{Eigenvector}).
    \end{itemize}
    
    \item \textbf{Extração de Atributos Intrínsecos:} Além das métricas estruturais, a pipeline extrai atributos já presentes nos dados, que descrevem características dos nós e arestas. Estes são classificados como:
    \begin{itemize}
        \item \textbf{Atributos Categóricos:} Rótulos discretos, como \texttt{node\_label}.
        \item \textbf{Atributos Contínuos:} Vetores de características numéricas, como \texttt{node\_attr\_*}.
    \end{itemize}
    
    \item \textbf{Armazenamento Tabular:} Todos os atributos extraídos, tanto de nós quanto de arestas, são consolidados em tabelas de dados (arquivos CSV). Cada linha nessas tabelas corresponde a um nó (ou aresta) e contém seu identificador de grafo (\texttt{graph\_id}), o rótulo da classe do grafo (\texttt{graph\_label}) e os valores brutos de todos os atributos calculados e extraídos. Este formato tabular serve como a entrada padronizada para a Etapa Q.
\end{enumerate}

Essa etapa é crucial, pois transforma o problema de classificação de estruturas de tamanho variável em um problema de classificação de vetores de características, embora ainda em um formato intermediário (por nó/aresta) que será agregado na etapa seguinte.

\subsection{Etapa Q: Geração da Representação Binária}

A etapa de quantização, implementada em \texttt{pipeline.py}, é o coração da nossa metodologia. Ela converte as tabelas de atributos geradas na Etapa M em um único vetor binário de tamanho fixo para cada grafo. O processo distingue entre atributos contínuos e categóricos.

\subsubsection{Atributos Contínuos: KernelCanvas++}

Para atributos contínuos, como as métricas de centralidade, utilizamos uma abordagem baseada em \textit{KernelCanvas++}. Para cada atributo, um conjunto de $n_{kernels}$ é posicionado no espaço de valores do atributo. Esses kernels funcionam como pontos de referência. A representação binária de um valor de atributo é gerada ativando os kernels mais próximos a ele.

A seleção dos centros dos kernels é crucial. A pipeline suporta duas estratégias:
\begin{enumerate}
    \item \textbf{K-Means}: Uma abordagem não supervisionada onde os centros são os centróides de clusters dos valores observados no treino.
    \item \textbf{FPS+MI}: Uma abordagem supervisionada que utiliza \textit{Farthest Point Sampling} (FPS) para garantir uma boa cobertura do espaço de valores e, em seguida, seleciona os kernels que possuem maior Informação Mútua (MI) com os rótulos das classes, priorizando características mais discriminativas.
\end{enumerate}
Para cada grafo, os valores de um atributo contínuo são transformados em uma sequência de vetores binários, que são então agregados (e.g., por soma lógica OR bit a bit) para formar uma assinatura binária única para aquele atributo no grafo.

\subsubsection{Atributos Categóricos: Histograma-Termômetro}

Para atributos categóricos (e.g., \texttt{node\_label}), a abordagem é diferente. Para cada grafo, um histograma é construído, contando a ocorrência de cada valor de rótulo possível. Em seguida, a contagem de cada categoria é codificada usando um \textit{código de termômetro}.

A codificação de termômetro representa um valor inteiro (a contagem) como um vetor binário onde os primeiros `level` bits são 1 e o restante é 0. A pipeline implementa duas variantes:
\begin{itemize}
    \item \textbf{Termômetro Linear}: Os `levels` são linearmente espaçados entre a contagem mínima (0) e máxima.
    \item \textbf{Termômetro Distributivo}: Os `levels` são definidos pelos quantis da distribuição de contagens observada no conjunto de treinamento, tornando a codificação mais robusta a outliers.
\end{itemize}

\subsection{Alocação Automática de Orçamento de Bits}

A definição manual dos hiperparâmetros de binarização (e.g., $n_{kernels}$, $bits\_per\_kernel$) para cada atributo é impraticável. O nosso sistema implementa uma alocação automática de orçamento de bits.

O processo, definido em \texttt{\_plan\_budget\_allocation}, segue estes passos:
\begin{enumerate}
    \item Um orçamento total de bits, $B$, é definido para a representação final do grafo.
    \item Este orçamento é dividido entre atributos \textit{estruturais} (centralidades) e \textit{não-estruturais} (atributos brutos dos nós/arestas) usando um fator $\alpha$:
    \begin{equation}
        B_s = \alpha \cdot B \quad ; \quad B_{ns} = (1 - \alpha) \cdot B
    \end{equation}
    \item Para cada atributo $i$, um orçamento individual $\beta_i$ é sorteado aleatoriamente do pool correspondente ($B_s$ ou $B_{ns}$), garantindo um mínimo de bits (\texttt{min\_info}).
    \item O orçamento $\beta_i$ é então fatorado aleatoriamente em $n_{kernels}$ e $bits\_per\_kernel$ (para atributos contínuos) ou usado para definir o número de bits do termômetro (para categóricos).
\end{enumerate}
Este procedimento automatiza a configuração, introduz diversidade nas representações e permite uma exploração eficiente do espaço de hiperparâmetros.

\section{Classificação com WiSARD}

O vetor binário final de cada grafo, resultante da concatenação das representações de todos os seus atributos, é então usado para treinar um classificador WiSARD. A WiSARD é composta por um conjunto de \textit{discriminadores}, um para cada classe do problema. Cada discriminador é uma coleção de RAMs.

Durante o treinamento, o vetor de entrada de um exemplo da classe $C$ é dividido em tuplas de $n$ bits. Cada tupla forma um endereço que aponta para uma posição em uma RAM do discriminador de $C$. O conteúdo dessa posição é então definido como 1. Na fase de teste, um novo vetor de entrada é apresentado a todos os discriminadores. A resposta de cada discriminador é a soma dos conteúdos das posições de memória endereçadas. O discriminador com a maior soma "reconhece" o padrão e sua classe é atribuída à entrada.

\section{Avaliação Experimental}

Para validar a eficácia da pipeline M+Q, conduzimos uma série de experimentos em datasets de benchmark da coleção TUDataset \cite{morris2020}. O objetivo foi avaliar a qualidade da representação binária gerada pela pipeline, utilizando classificadores padrão como ferramenta de medição. A análise dos resultados, cujos dados brutos são gerenciados pelo script \texttt{analyze\_results.py}, é detalhada a seguir.

\subsection{Configuração Experimental}

Os experimentos foram executados sobre sete datasets da coleção TUDataset: \texttt{MUTAG}, \texttt{PROTEINS}, \texttt{NCI1}, \texttt{PTC\_MR}, \texttt{AIDS}, \texttt{ENZYMES} e \texttt{IMDB-MULTI}. O arquivo \texttt{experiments.json} contém 35 registros, correspondentes a cinco execuções da pipeline (uma por configuração de representação) e aos sete datasets processados em cada execução. O número de grafos varia de 188, no \texttt{PTC\_MR}, a 2.000, no \texttt{AIDS}; o número de classes varia de duas a seis.

No arquivo de resultados, as configurações efetivamente registradas foram:
\begin{itemize}
    \item \textbf{Configuração C1}: $n_{kernels}=8$, $bits\_per\_kernel=4$ e $k_{activate}=2$;
    \item \textbf{Configuração C2}: $n_{kernels}=16$, $bits\_per\_kernel=4$ e $k_{activate}=3$.
\end{itemize}

Em ambas, a estratégia de seleção de kernels registrada foi nula, o que corresponde ao caminho padrão baseado em ECDF e K-means no código da pipeline. Assim, embora o programa também implemente estratégias alternativas, como FPS, elas não foram utilizadas nos resultados analisados neste artigo. O parâmetro $k_{activate}$ define quantos kernels são ativados na transformação de cada valor contínuo. A dimensionalidade final depende do dataset e da configuração, variando de 64 a 1.280 bits nos registros válidos.

Para avaliar o poder discriminativo dos vetores binários resultantes, o modo \texttt{benchmark} executou cinco classificadores: Random Forest (RF), Support Vector Machine (SVM), k-Nearest Neighbors (kNN), Multi-Layer Perceptron (MLP) e WiSARD. O desempenho foi medido por validação cruzada estratificada de 10 folds, reportando a acurácia média, o desvio padrão, o F1 médio e a matriz de confusão agregada. Essa comparação separa a qualidade da representação da capacidade de um único classificador e permite verificar especificamente o comportamento do WiSARD.

\subsection{Resultados e Discussão}

Os melhores resultados de acurácia média por combinação de dataset e classificador estão sumarizados na Tabela~\ref{tab:results}. Para cada célula, foi selecionada a melhor entre C1 e C2; portanto, a tabela não representa uma única configuração global, mas o melhor resultado observado no pequeno grid experimental.

\begin{table}[ht]
\centering
\caption{Melhor acurácia média (\%) e desvio padrão entre C1 e C2 para cada dataset e classificador.}
\label{tab:results}
\resizebox{\textwidth}{!}{%
\begin{tabular}{lccccc}
\toprule
\textbf{Dataset} & \textbf{RF} & \textbf{SVM} & \textbf{kNN} & \textbf{MLP} & \textbf{WiSARD} \\
\midrule
AIDS & 99.65 $\pm$ 0.63 & 99.50 $\pm$ 0.71 & 99.45 $\pm$ 0.61 & 99.55 $\pm$ 0.65 & 80.00 $\pm$ 0.00 \\
MUTAG & 87.81 $\pm$ 5.74 & 86.75 $\pm$ 6.28 & 88.30 $\pm$ 3.99 & 89.42 $\pm$ 7.02 & 80.94 $\pm$ 6.60 \\
PROTEINS & 74.76 $\pm$ 5.44 & \textbf{76.11 $\pm$ 4.07} & 73.68 $\pm$ 4.47 & 73.77 $\pm$ 5.26 & 63.43 $\pm$ 3.01 \\
NCI1 & 67.62 $\pm$ 1.43 & 67.15 $\pm$ 1.94 & 66.89 $\pm$ 1.45 & 67.40 $\pm$ 1.17 & 57.74 $\pm$ 2.52 \\
PTC\_MR & \textbf{60.78 $\pm$ 5.63} & 56.39 $\pm$ 6.13 & 59.10 $\pm$ 10.07 & 58.74 $\pm$ 4.81 & 56.40 $\pm$ 1.72 \\
ENZYMES & \textbf{67.00 $\pm$ 5.15} & 55.00 $\pm$ 7.15 & 57.83 $\pm$ 3.66 & 60.67 $\pm$ 3.67 & 22.17 $\pm$ 3.80 \\
IMDB-MULTI & \textbf{48.13 $\pm$ 3.24} & 46.93 $\pm$ 2.94 & 39.73 $\pm$ 4.78 & 47.47 $\pm$ 3.12 & 45.73 $\pm$ 2.43 \\
\bottomrule
\end{tabular}
}
\end{table}

Os resultados evidenciam três padrões. Primeiro, \texttt{AIDS} é o dataset mais separável: todos os classificadores convencionais ultrapassam 99\%, enquanto o WiSARD alcança 80\%. Segundo, \texttt{MUTAG} também apresenta desempenho elevado, com o MLP atingindo 89,42\% e kNN 88,30\%. Terceiro, \texttt{ENZYMES}, \texttt{IMDB-MULTI} e \texttt{PTC\_MR} permanecem mais difíceis, com melhores resultados de 67,00\%, 48,13\% e 60,78\%, respectivamente.

Considerando a média dos melhores resultados por dataset, RF obteve 72,25\%, seguido de MLP (71,00\%), SVM (69,69\%), kNN (69,28\%) e WiSARD (58,06\%). O RF foi o melhor classificador em quatro datasets; SVM foi superior em \texttt{PROTEINS}, MLP em \texttt{MUTAG} e RF voltou a liderar em \texttt{AIDS}, \texttt{ENZYMES}, \texttt{NCI1}, \texttt{PTC\_MR} e \texttt{IMDB-MULTI}. A vantagem do RF sobre o SVM em \texttt{ENZYMES} foi de 12 pontos percentuais, sugerindo que a representação agregada contém padrões não lineares ou interações entre atributos que são melhor explorados por conjuntos de árvores.

Uma análise diagnóstica importante, implementada em \texttt{pipeline.py}, é a correlação entre a distância no espaço de valores original e a distância de Hamming no espaço binário. No \texttt{AIDS}, o atributo \texttt{node\_attr\_1} apresentou correlação de 0,781, o maior valor observado entre os atributos desse dataset; \texttt{degree} e \texttt{pagerank} apresentaram 0,603 e 0,451, respectivamente. Esses valores indicam que a codificação preserva parcialmente a ordem métrica dos dados, mas não de forma uniforme para todos os atributos. No \texttt{ENZYMES}, por exemplo, as correlações variaram de 0,145 a 0,586, o que é compatível com a maior dificuldade de classificação desse conjunto.

Também há uma diferença relevante entre aumentar a resolução da representação e melhorar a acurácia. A configuração C2, com 16 kernels, foi a melhor para a maioria dos classificadores e datasets, mas C1 permaneceu competitiva em casos como \texttt{AIDS}, \texttt{MUTAG} e \texttt{PTC\_MR}. Portanto, o aumento do vetor não garante ganho monotônico: a representação precisa ser avaliada junto ao classificador e ao conjunto de grafos. Além disso, como o arquivo agregado contém execuções de configurações distintas e não repetições independentes de uma mesma configuração, os desvios padrão reportados são os obtidos nos folds da validação cruzada, e não uma estimativa de variabilidade entre execuções completas.

\section{Conclusão}

Este trabalho apresentou a M+Q, uma pipeline integrada e automatizada para a classificação de grafos com Redes Neurais sem Peso. A metodologia aborda de ponta a ponta o processo de extração de atributos, binarização e otimização de hiperparâmetros, culminando em uma representação binária robusta para o classificador WiSARD.

A combinação de KernelCanvas++ para atributos contínuos, codificação de histograma-termômetro para dados categóricos e, crucialmente, a alocação automática de orçamento de bits, provou ser uma estratégia eficaz. Os resultados experimentais indicam que a supervisão na etapa de binarização (via \texttt{fps+mi}) e a automação da alocação de recursos de representação são chaves para alcançar alta acurácia.

Como trabalhos futuros, pretendemos explorar sistematicamente as estratégias FPS e FPS+MI, comparar a alocação automática de orçamento de bits com o grid fixo utilizado nos resultados atuais e investigar funções de agregação mais sofisticadas. Também será necessário repetir as configurações com sementes controladas e avaliar a WiSARD com uma seleção de hiperparâmetros própria, pois o benchmark atual mostra que a representação já é útil para classificadores convencionais, mas ainda não produz desempenho equivalente no classificador-alvo em todos os datasets.

\bibliographystyle{sbc}
\bibliography{sbc-template} % Assumindo que existe um arquivo sbc-template.bib

\end{document}
